\section{Related Work}
Articulated transport provides the mechanical baseline, cable-connected systems motivate mode strategies, and hybrid optimization provides tools for their numerical treatment.

\subsection{Articulated Vehicles and Warehouse Automation}
Articulated planning couples nonholonomic motion, multiple body footprints, and successive headings. Ljungqvist et al.~\cite{ljungqvist2019} combine planning and tracking for a general two-trailer vehicle. Deniz and Auat Cheein~\cite{deniz2025} develop predictive path-following for generalized trailer trains with dynamic obstacles. Tracailer~\cite{xu2025} integrates a compact smooth trajectory representation with efficient initial-path search. GATRO~\cite{gatro2026} exploits semidefinite structure and graphics processing unit (GPU) acceleration. These methods advance rigid-coupling transport; flexible connections introduce additional relative configurations and switching behavior.

Broader warehouse and mobile-robot studies address complementary issues. Bidirectional--unidirectional sampling~\cite{wang2022rrt} improves difficult kinodynamic connections, while dense-environment coordination~\cite{xu2024dense} couples task assignment with motion. Learned parking trajectories~\cite{chai2023} and two-stage coverage optimization~\cite{painting2025} illustrate the role of initial solutions and discrete decisions. Universal differential-drive planning~\cite{ddr2025} and spatial--temporal planning for car-like robots~\cite{han2024} emphasize representations matched to platform constraints. Our numerical study fixes those constraints and compares cable-mode handling.

\subsection{Cable-Connected Transport}
Xiao et al.~\cite{xiao2021} model leash-mediated hybrid interaction in a guide-dog system. Yang et al.~\cite{yang2022} coordinate quadrupeds towing a common load and optimize alternative cable-mode combinations. Xu et al.~\cite{xu2024cable} use observers and distributed model predictive control (MPC) for quadrotor--quadruped towing. The closest mechanical precedent is Zhang et al.~\cite{zhang2026}, where a quadruped tows a wheeled trailer and a searched mode sequence guides trajectory optimization. Its dynamic plant differs from a car-like tractor, but the strategy of selecting modes before continuous refinement is relevant.

Tether geometry introduces another distinction. Mart\'inez-Rozas et al.~\cite{marsupial2023} plan coupled trajectories for an unmanned aerial vehicle (UAV) and unmanned ground vehicle (UGV), including a non-taut tether with adjustable length. Here the material cable length is constant. Cable bending accommodates shorter endpoint distances as slack. For serial towing, each passive cart transmits motion to the next through its own rear attachment.

\subsection{Numerical Strategies and Body Geometry}
Contact-implicit MPC~\cite{cleach2024} and consensus complementarity control~\cite{aydinoglu2024} use structured treatments of contact-mode constraints. Hybrid iterative linear quadratic regulation (iLQR)~\cite{kong2023} uses mode-sensitive derivatives; local smoothing~\cite{pang2023} and bilevel optimization~\cite{shirai2024} offer alternative treatments of switching. These works motivate initialization, complementarity, and mode-selection strategies for the present nonlinear program (NLP). A mixed-integer nonlinear program (MINLP) provides an explicit binary representation, while log-sum-exp (LSE) offers an optional approximation of a maximum.

Geometry and trajectory representation remain separate design choices. Geometrically constrained multicopter optimization~\cite{wang2022minco} develops an efficient trajectory representation, and graphs of convex sets~\cite{kurtz2023} combine route choices with continuous optimization. Shape-changing robots~\cite{shape2023} and implicit signed-distance planning~\cite{anyshape2023} demonstrate the importance of complete body geometry. We retain the OBCA distance certificate~\cite{obca2021} for every optimized body--obstacle, body--body, and wall constraint.

The comparison extracts the fixed-mode strategy from~\cite{zhang2026} and the alternative-mode idea from~\cite{yang2022}, implements them on the common cable model, and separately tests equivalent mathematical representations. The resulting experiments evaluate the transferred strategies under shared kinematics, geometry, and computational conditions.
